\section{Every coefficient sheaf on the current source: exact gluing and
supported
cohomology}\label{every-coefficient-sheaf-on-the-current-source-exact-gluing-and-supported-cohomology}

24 September 2026. Independent continuation of
\url{SOURCE_SHEAF_SPECIALIZATION.md}, SS0--SS9. The coefficient
comparison left in SS7 is now an explicit object with inverse
construction, all morphisms, its complete restriction defect and its
derived supported complex.

\subsection{CGS0. Source and operation
prerequisites}\label{cgs0.-source-and-operation-prerequisites}

The current \nolinkurl{CORPUS_AND_OPERATION_RULES.md} (private source record), including the
global-argument and classical-boundary corrections, was read for this
continuation. The source remains the current B1--B5, with the verbatim
U01--U18 and whole-history arguments retained as in ST0, AC0 and SS0. In
particular \(Z_1/\tau\) has no source addition, coordinate or \(Z_2\)
parity. The complete reconstructed integer layer is already present
before any arithmetic receiver below is used. Nothing here changes those
definitions or treats a coefficient choice as a counterexample to the
user's complete programme.

The proved source receiving space is \[
 X=U\cup\{\mathfrak m\},\qquad U=X_{\mathbb Z}\simeq\operatorname{Spec}\mathbb Z,
\] where the only open neighbourhood of \(\mathfrak m\) is \(X\). Keep
\[
 \mathcal R=j_*\mathcal O_U\times i_*\mathbb Z,\qquad
 R_0=\mathcal R(X)=\mathbb Z[\epsilon]/(\epsilon^2-\epsilon)
 \cong\mathbb Z\times\mathbb Z .
\tag{CGS0.1}
\] Here \(\epsilon=[1]\) is the receiving idempotent, not primitive
\(\tau\). The original source integer \(n\) maps to \(n\epsilon\). The
bracketed image of \(\tau\) is the formal global unit of this receiver;
its new additive operations are not pulled back to the source. All
sheaves and complexes in this document have these explicitly constructed
receiving modules as coefficients.

The letter \(i\) in the sheaf-support calculations denotes the closed
topological inclusion equipped, when module rings are needed, with the
full stalk ring \(R_0=i^{-1}\mathcal R\). This matters: the selected
residue ring \(\mathbb Z\) through second projection is a different
coefficient functor, recorded in CGS7. We do not replace the full stalk
by that residue.

\subsection{CGS1. A complete classification, including inverse
functors}\label{cgs1.-a-complete-classification-including-inverse-functors}

The human-source gluing mechanism is explicitly recalled by Pierre
Deligne in \emph{La conjecture de Weil. II}, §3.4.8, printed p.209: an
open sheaf, a closed sheaf, and its specialization map into the
restriction of the open direct image describe the whole sheaf. The
calculation here specializes that recalled equivalence to the proved
ringed space CGS0.1 and separates its two receiving idempotents. It is
not presented as a new discovery of the general sheaf-gluing
equivalence. CGS10 gives the exact specialization target and source
record.

Define a gluing datum to consist of \[
 (\mathcal G,A_+,A_-,r),\qquad
 r:A_+\longrightarrow\Gamma(U,\mathcal G),
\tag{CGS1.1}
\] where \(\mathcal G\) is a sheaf of \(\mathcal O_U\)-modules, \(A_+\)
and \(A_-\) are abelian groups, and \(r\) is \(\mathbb Z\)-linear. No
surjectivity, injectivity, topology on these groups or purity is
imposed.

Every sheaf \(\mathcal M\) of \(\mathcal R\)-modules supplies exactly
such data: \[
 \mathcal G=j^{-1}\mathcal M,\qquad
 A=\mathcal M(X)=\mathcal M_{\mathfrak m},\qquad
 A_+=\epsilon A,\quad A_-=(1-\epsilon)A.
\tag{CGS1.2}
\] The restriction \(A\to\Gamma(U,\mathcal G)\) kills \(A_-\), because
\(1-\epsilon\) restricts to zero on \(U\). Its restriction to \(A_+\) is
\(r\). The decomposition \(A=A_+\oplus A_-\) is proved by the
complementary idempotent identities, as in ST5.

Conversely construct a sheaf \(\mathcal M(\mathcal G,A_+,A_-,r)\) as
follows: \[
 \mathcal M(X)=A_+\oplus A_-,\qquad
 \mathcal M(V)=\mathcal G(V)\quad(V\subseteq U),
\tag{CGS1.3}
\] and \[
 \operatorname{res}_{X,V}(a_+,a_-)=r(a_+)|_V.
\tag{CGS1.4}
\] The \(R_0\)-module action on \(A_+\oplus A_-\) is componentwise:
under \(R_0\cong\mathbb Z\times\mathbb Z\),
\((m,n)(a_+,a_-)=(ma_+,na_-)\). On proper opens it is the given
\(\mathcal O_U\)-module action. The restrictions are module maps because
\(R_0\to\mathcal O_U(V)\) is first projection followed by the integer
localization, and \(r\) is \(\mathbb Z\)-linear.

This is a sheaf. A cover of \(X\) contains \(X\) itself, so its gluing
condition is automatic from the matching condition. A cover of a proper
open is a cover inside \(U\), where \(\mathcal G\) is already a sheaf.
These are all covers. Applying CGS1.2 to this construction recovers the
given quadruple. Applying CGS1.3 to an original sheaf recovers all its
sections and all restrictions. These are explicit inverse functors, not
merely a correspondence between dimensions.

A morphism between two data is precisely \[
 (g,a_+,a_-),\quad
 g:\mathcal G\to\mathcal G',\quad
 a_\pm:A_\pm\to A'_\pm,\quad
 (\Gamma(U,g))\,r=r'a_+.
\tag{CGS1.5}
\] It acts by \(a_+\oplus a_-\) on \(X\) and by \(g\) on proper opens.
The displayed equation is exactly compatibility with restriction. Every
module-sheaf morphism has this form because it commutes with
\(\epsilon\). Composition is componentwise and preserves the displayed
equation. This proves an equivalence of categories in both directions,
including all maps.

\subsection{CGS2. Exactness, kernels and
cokernels}\label{cgs2.-exactness-kernels-and-cokernels}

An exact sequence of these sheaves is exact on their \(A_+\), \(A_-\),
and open-sheaf components. Necessity follows by taking the stalk at
\(\mathfrak m\), then the two idempotents, and taking all stalks in
\(U\). These conditions are sufficient because they exhaust the stalks
of \(X\).

For a morphism CGS1.5, its kernel has data \[
 \left(\ker g,\ \ker a_+,\ \ker a_-,\ r|_{\ker a_+}\right).
\tag{CGS2.1}
\] The last map takes values in \(\Gamma(U,\ker g)\), because
\((\Gamma g)r=r'a_+\) vanishes on \(\ker a_+\), and a section in the
kernel is exactly a section of the kernel sheaf.

Its cokernel has open sheaf \(\operatorname{coker}g\) and stalk groups
\(\operatorname{coker}a_\pm\). The restriction map is induced by \[
 A'_+\xrightarrow{r'}\Gamma(U,\mathcal G')
 \longrightarrow\Gamma(U,\operatorname{coker}g).
\tag{CGS2.2}
\] It kills \(\operatorname{im}a_+\) by the compatibility equation,
hence descends. This is the cokernel datum by stalkwise verification. In
particular \(\Gamma(U,\operatorname{coker}g)\) is not replaced by the
cokernel of global sections unless that additional equality has actually
been proved.

The same description gives images, finite limits and finite sums. It
proves directly that the classified category has the abelian structure
of the original module-sheaf category. No abelian category of primitive
\(\tau\) objects is assumed.

\subsection{CGS3. The exact open and closed
objects}\label{cgs3.-the-exact-open-and-closed-objects}

The three standard sheaves have these complete data: \[
 \begin{array}{c|c}
 \text{sheaf}&(\mathcal G,A_+,A_-,r)\\ \hline
 j_!\mathcal G&(\mathcal G,0,0,0)\\
 j_*\mathcal G&(\mathcal G,\Gamma(U,\mathcal G),0,\mathrm{id})\\
 i_*A&(0,\epsilon A,(1-\epsilon)A,0).
 \end{array}
\tag{CGS3.1}
\] For \(j_!\), these formulas also follow directly from extension by
zero: a nonzero section supported away from \(\mathfrak m\) cannot
extend by zero globally, since the closure in \(X\) of every nonempty
subset of \(U\) contains \(\mathfrak m\). Its value on \(X\) is
therefore zero. It is the original sheaf on every proper open. The
formulas for \(j_*\) and \(i_*\) follow from their actual section
definitions.

For a general \(\mathcal M\), the restriction comparison \[
 \eta_{\mathcal M}:\mathcal M\longrightarrow j_*\mathcal G
\tag{CGS3.2}
\] is identity on \(U\), and on \(X\) it is \((a_+,a_-)\mapsto r(a_+)\).
Thus its \emph{complete} support kernel and cokernel are \[
 \ker\eta_{\mathcal M}
 =i_*(\ker r\oplus A_-),\qquad
 \operatorname{coker}\eta_{\mathcal M}=i_*(\operatorname{coker}r).
\tag{CGS3.3}
\] In the first kernel, \(\epsilon\) acts as identity on \(\ker r\) and
as zero on \(A_-\); in the cokernel it acts as identity. This retains
the two source-idempotent types inside the same closed support. The
exact sequence is \[
 0\to i_*(\ker r\oplus A_-)\to\mathcal M
 \xrightarrow{\eta_{\mathcal M}}j_*\mathcal G
 \to i_*(\operatorname{coker}r)\to0.
\tag{CGS3.4}
\]

The source's automatic boundary component is \(i_*A_-\), and it splits
off: \[
 \mathcal M\cong
 \mathcal M(\mathcal G,A_+,0,r)\oplus i_*A_-.
\tag{CGS3.5}
\] The explicit maps are the two idempotent projections and their
inclusions. The additional supported objects \(\ker r\) and
\(\operatorname{coker}r\) belong to the \(\epsilon=1\) component.
Consequently the automatic \(\epsilon=0\) splitting does not erase them.

\subsection{\texorpdfstring{CGS4. The gluing extension is exactly the
map
\(r\)}{CGS4. The gluing extension is exactly the map r}}\label{cgs4.-the-gluing-extension-is-exactly-the-map-r}

There is another exact sequence whose maps are also fully determined: \[
 0\longrightarrow j_!\mathcal G
 \longrightarrow\mathcal M
 \longrightarrow i_*(A_+\oplus A_-)
 \longrightarrow0.
\tag{CGS4.1}
\] On \(U\), its first map is identity and its last target is zero. At
\(\mathfrak m\), its first stalk is zero and its last map is identity.
This proves exactness. The last map on global sections is also identity;
no restriction lift is hidden there.

The sequence is the actual pullback of \[
 0\to j_!\mathcal G\to j_*\mathcal G
 \to i_*\Gamma(U,\mathcal G)\to0
\tag{CGS4.2}
\] along \(i_*r:i_*(A_+\oplus A_-)\to i_*\Gamma(U,\mathcal G)\), where
\(r\) kills \(A_-\). Explicitly, \[
 \mathcal M=
 j_*\mathcal G\times_{\,i_*\Gamma(U,\mathcal G)}
 i_*(A_+\oplus A_-).
\tag{CGS4.3}
\] On \(X\) the pairs are \((g,a)\) with \(g=r(a_+)\), hence are exactly
\(A_+\oplus A_-\). On a proper open they are \(\mathcal G(V)\). These
identifications preserve every module action and restriction, proving
the formula as a sheaf isomorphism.

The complete classification of extensions with these fixed endpoints is
\[
 \boxed{
 \operatorname{Ext}_{\mathcal R}^{1}(i_*A,j_!\mathcal G)
 \cong\operatorname{Hom}_{\mathbb Z}(\epsilon A,\Gamma(U,\mathcal G)).
 }
\tag{CGS4.4}
\] Indeed in any such extension, its open sheaf is canonically
\(\mathcal G\) and its closed stalk is canonically \(A\), by exactness.
The only remaining datum in CGS1 is \(r\), which must kill
\((1-\epsilon)A\). Conversely CGS4.3 constructs the extension for every
such map. An equivalence of extensions fixing the endpoints fixes those
same open and stalk identifications, hence fixes \(r\). Finally the Baer
sum is addition of these maps: form the fibre product over \(i_*A\),
then push out its two open kernels along addition
\(\mathcal G\oplus\mathcal G\to\mathcal G\). The resulting restriction
on \(A\) is \(r_1+r_2\). This proves the group isomorphism, not just a
set correspondence.

In particular CGS4.1 has a sheaf section precisely when \(r=0\). A
section would be identity on the \(A\)-stalk and zero on the open
source; compatibility with restriction forces \(r=0\). Conversely
\(r=0\) gives the direct sum \(j_!\mathcal G\oplus i_*A\) and the
indicated section.

This is the gluing extension, not the support/restriction sequence
CGS3.4. The exact pullback CGS4.3 displays their relationship and
prevents their different lifting questions from being conflated.

\subsection{CGS5. Higher supported cohomology for every coefficient
sheaf}\label{cgs5.-higher-supported-cohomology-for-every-coefficient-sheaf}

The full global-section functor on \(X\) is exact: \[
 \Gamma(X,\mathcal M)=\mathcal M_{\mathfrak m}=A_+\oplus A_-.
\tag{CGS5.1}
\] The only neighbourhood of \(\mathfrak m\) is \(X\), so this is
equality with a stalk functor, which is exact. Thus
\(H^q(X,\mathcal M)=0\) for \(q>0\), on this stated site.

For \(Z=\{\mathfrak m\}\), form the supported-section kernel \[
 \Gamma_Z(X,\mathcal M)=
 \ker\bigl(\Gamma(X,\mathcal M)\to\Gamma(U,\mathcal G)\bigr).
\] Its restriction triangle and CGS5.1 give, with \textbf{all} higher
open cohomology retained, \[
 \boxed{
 \begin{aligned}
 H_Z^0(X,\mathcal M)&=A_-\oplus\ker r,\\
 H_Z^1(X,\mathcal M)&=\operatorname{coker}r,\\
 H_Z^q(X,\mathcal M)&=H^{q-1}(U,\mathcal G)\quad(q\ge2).
 \end{aligned}}
\tag{CGS5.2}
\] To derive these formulas directly, take a flabby resolution
\(\mathcal G\to\mathcal I^\bullet\) on \(U\), put
\(B^\bullet=\Gamma(U,\mathcal I^\bullet)\), and embed \(r\) in its
degree-zero cocycles through
\(\Gamma(U,\mathcal G)\hookrightarrow B^0\). The complete supported
complex is \[
 K^0=A_+\oplus A_-,\qquad
 K^q=B^{q-1}\ (q\ge1),
\] \[
 d_K^0(a_+,a_-)=r(a_+),\qquad
 d_K^q=-d_B^{\,q-1}\ (q\ge1).
\tag{CGS5.3}
\] The differential squares to zero because \(r(a_+)\) lies in the
cocycles of \(B^0\). Its degree-zero kernel is the first line of CGS5.2.
Its degree-one cocycles are \(\Gamma(U,\mathcal G)\) and its degree-one
boundaries are \(r(A_+)\), proving the second line. Its remaining
cohomology is the displayed shifted open cohomology. This also fixes
every sign.

The long exact sequence's degree-zero connecting map is precisely \[
 \Gamma(U,\mathcal G)\longrightarrow\operatorname{coker}r,
 \qquad g\longmapsto[g].
\tag{CGS5.4}
\] For \(q\ge1\), its connecting map \[
 H^q(U,\mathcal G)\longrightarrow H_Z^{q+1}(X,\mathcal M)
\] is the isomorphism given by including a cocycle \(b\in B^q\) as the
same \(b\) in \(K^{q+1}\), with the sign convention CGS5.3. These are
actual maps, not unspecified groups that might contain an obstruction.

The local cohomology sheaves are the skyscrapers with these groups: \[
 \mathcal H_Z^q(\mathcal M)=i_*H_Z^q(X,\mathcal M).
\tag{CGS5.5}
\] Away from \(Z\), every such stalk vanishes. At \(\mathfrak m\), the
only neighbourhood is \(X\), giving the displayed group. Its
\(R_0\)-action is the one already stated: \(A_-\) has \(\epsilon=0\),
whereas \(\ker r\), \(\operatorname{coker}r\) and every higher
open-cohomology term have \(\epsilon=1\). This is an explicit action
calculation, not a numerical purity assignment.

For completeness, the derived open direct image has \[
 R^0j_*\mathcal G=j_*\mathcal G,\qquad
 R^qj_*\mathcal G=i_*H^q(U,\mathcal G)\quad(q>0).
\tag{CGS5.6}
\] At \(\mathfrak m\) this follows from the same sole neighbourhood. On
\(U\), the higher sheaves vanish because the restricted direct-image
functor is identity; equivalently, a sheaf-cohomology class becomes zero
in a sufficiently small neighbourhood of its representative stalk. Thus
these formulas include every possible higher term.

\subsection{CGS6. Complexes: the full mapping fibre and its boundary
representatives}\label{cgs6.-complexes-the-full-mapping-fibre-and-its-boundary-representatives}

For a cochain complex \(\mathcal M^\bullet\) of \(\mathcal R\)-modules,
the classification is termwise exact: \[
 (\mathcal G^\bullet,A_+^\bullet,A_-^\bullet,r^\bullet),
\qquad
 r^\bullet:A_+^\bullet\to\Gamma(U,\mathcal G^\bullet),
\tag{CGS6.1}
\] where the differentials are the classified sheaf morphisms. Thus
\(r^\bullet\) is a chain map, and conversely these data construct a
complex by CGS1. Chain maps and homotopies are classified by the same
compatibility equation in each degree. This assertion applies without
any boundedness restriction.

Here is a fully specified derived calculation for bounded-below
complexes, including all two-term and finite complexes used in the
programme. Take a functorial flabby double resolution of
\(\mathcal G^\bullet\) and its total complex; in each degree only
finitely many summands occur because the original complex is bounded
below and the resolution degrees are nonnegative. Let \[
 B^\bullet=R\Gamma(U,\mathcal G^\bullet)
\] denote that explicit total section complex, and let \[
 \widetilde r:A_+^\bullet\longrightarrow B^\bullet
\tag{CGS6.2}
\] be the composite of \(r^\bullet\) with the resolution map. A usual
double resolution is obtained degree by degree from injective or
Godement resolutions and their functorial maps; its resolution-direction
signs are included in the total differential. No higher open cohomology
is deleted.

The actual supported complex is the following mapping fibre: \[
 \boxed{
 K^n=A_+^n\oplus A_-^n\oplus B^{n-1},\quad
 d_K(a_+,a_-,b)
 =(d_+a_+,d_-a_-,\widetilde r(a_+)-d_Bb).
 }
\tag{CGS6.3}
\] To check \(d_K^2=0\), its first two components are the squares of the
original differentials, and its last is
\(\widetilde r(d_+a_+)-d_B\widetilde r(a_+)+d_B^2b=0\). This convention
is a concrete model of \[
 R\Gamma_Z(X,\mathcal M^\bullet)
 \cong A_-^\bullet\oplus
 \operatorname{Fib}\!\left(A_+^\bullet
       \xrightarrow{\widetilde r}R\Gamma(U,\mathcal G^\bullet)\right).
\tag{CGS6.4}
\] The restriction triangle proves the identification: the middle global
derived sections equal \(A_+^\bullet\oplus A_-^\bullet\) because CGS5.1
is exact, the open term is \(B^\bullet\), and its restriction is
\((\widetilde r,0)\). The mapping fibre of precisely that map is CGS6.3.
The sheaf-level supported complex is \(i_*K^\bullet\).

For genuinely unbounded input, the termwise classification CGS6.1 is
unchanged. The same mapping-fibre formula applies when
\(R\Gamma(U,\mathcal G^\bullet)\) is represented by a K-injective model
of that unbounded complex; it must not be replaced by the
finite-diagonal totalization just used. The construction and
applications proved here require only the bounded-below case. No
conclusion about an unbounded original programme object is obtained by
truncating it.

Every connecting map is now explicit. If \(b\in B^n\) is a cocycle, its
boundary is \[
 \partial[b]=[(0,0,b)]\in H^{n+1}(K).
\tag{CGS6.5}
\] If \(b=d_Bc\), this representative is \(d_K(0,0,-c)\), so the formula
is well-defined. The class \([b]\) has a global lift exactly when there
exist a cocycle \(a_+\in A_+^n\) and \(c\in B^{n-1}\) with \[
 b=\widetilde r(a_+)-d_Bc.
\tag{CGS6.6}
\] This is not an assumed existence statement: it is the explicit
boundary criterion obtained by writing out when the representative in
CGS6.5 is a boundary. The already present \(A_-^n\) supplies an
independent global cocycle component and does not alter the restriction.

In every degree there is the complete exact sequence \[
 \begin{aligned}
 0\longrightarrow&
 \operatorname{coker}\!\left(
 H^{n-1}(A_+^\bullet)\to H^{n-1}(B^\bullet)\right)
 \longrightarrow H^n(K^\bullet)\\
 \longrightarrow&
 \ker\!\left(H^n(A_+^\bullet)\to H^n(B^\bullet)\right)
 \oplus H^n(A_-^\bullet)\longrightarrow0 .
 \end{aligned}
\tag{CGS6.7}
\] It follows by the long exact cohomology sequence of the displayed
mapping fibre, or directly by taking the first two coordinates of a
cocycle in CGS6.3. The \(H^n(A_-^\bullet)\) summand splits canonically
at the complex level; the remaining short exact sequence is not declared
split without a constructed section.

\subsection{CGS7. Full support and the selected residue are different
computed
functors}\label{cgs7.-full-support-and-the-selected-residue-are-different-computed-functors}

The full closed-point coefficient ring is
\(R_0=\mathbb Z\times\mathbb Z\), and \(i_*A\) in CGS3 allows either
idempotent type. Its supported right adjoint is computed by CGS5--CGS6.

If instead the selected closed receiving ring is \(\mathbb Z\) via the
second projection \(R_0\to\mathbb Z\), its pushforward has
\(\epsilon=0\). Its right adjoint on module sheaves is simply \[
 \mathcal M\longmapsto (1-\epsilon)\mathcal M_{\mathfrak m}=A_-.
\tag{CGS7.1}
\] Indeed a module map from a second-component stalk module into
\(\mathcal M\) lands in the \(\epsilon=0\) part of its stalk, which
automatically restricts to zero; conversely any such group map gives a
sheaf map. This proves the adjunction. The functor is exact because
stalk and idempotent projection are exact, so its derived version on
complexes is \(A_-^\bullet\).

It follows that the selected residue functor records the automatic
source boundary but not the entire topological-support functor. The
precise additional object is \[
 \operatorname{Fib}\!\left(A_+^\bullet
       \xrightarrow{\widetilde r}R\Gamma(U,\mathcal G^\bullet)\right),
\tag{CGS7.2}
\] with every group and map computed above. This is the explicit
relation between the two functors; no absence of a comparison is used to
dismiss their connection.

For the structure sheaf itself, CGS9 below shows that CGS7.2 is acyclic,
so the two support calculations agree there. They need not be silently
equated for other coefficients.

\subsection{CGS8. Commuting arithmetic actions and the gluing class in
higher
degrees}\label{cgs8.-commuting-arithmetic-actions-and-the-gluing-class-in-higher-degrees}

Suppose specified additive coefficient operators act over the identity
on the source space and commute with the supplied source action.
Commutation with integer \(1\) means commutation with \(\epsilon\). On a
localized arithmetic open an additive map commutes with every integer,
and therefore with every inverted integer: multiply the desired
commutation relation by that invertible multiplication operator to
verify it. Thus these maps are precisely the module maps in the
classification.

An action of a specified group or monoid \(\mathcal H\) consequently
consists exactly of compatible actions on \[
 \mathcal G,\quad A_+,\quad A_-,
\qquad
 \Gamma(U,F_h)\,r=rF_h^+\quad(h\in\mathcal H).
\tag{CGS8.1}
\] Their composition, inverses when specified, commutation relations and
chain-map properties hold on the sheaf exactly when they hold on these
entries. Conversely CGS1 constructs the whole sheaf action from those
data. A functorial resolution retains the actions in CGS6, so its cone
differential and every connecting map are equivariant. If an intended
geometric operator moves the base rather than fixing it, that requires
its stated pullback map; it is not silently included among these
over-identity actions.

In particular a finite-dimensional Frobenius-equivariant quotient of
\(\operatorname{coker}r\) has exactly the quotient action induced from
\(\Gamma(U,\mathcal G)\). A generalized eigenclass can have nonzero
boundary only at the same eigenvalue:
\((F-\alpha)^k\partial=\partial(F-\alpha)^k\) follows by the proved
equivariance. Thus CGS5.4 retains the complete receiving spectral
action; source-idempotent splitting alone does not change its
eigenvalue. No numerical weight of primitive \(\tau\) is introduced.

There is also a complete higher derived form of CGS4.4: \[
 R\operatorname{Hom}_{\mathcal R}(i_*A,j_!\mathcal G)
 \cong
 R\operatorname{Hom}_{R_0}
       \bigl(A,R\Gamma(U,\mathcal G)\bigr)[-1]
 \cong
 R\operatorname{Hom}_{\mathbb Z}
       \bigl(\epsilon A,R\Gamma(U,\mathcal G)\bigr)[-1].
\tag{CGS8.2}
\] Here the open derived sections carry \(\epsilon=1\). To prove it,
choose an injective resolution \(\mathcal I^\bullet\) of \(\mathcal G\).
For every resolution term the exact sequence \[
 0\to j_!\mathcal I^q\to j_*\mathcal I^q
 \to i_*\Gamma(U,\mathcal I^q)\to0
\] is CGS4.2. Hence it yields the triangle \[
 j_!\mathcal G\to Rj_*\mathcal G
 \to i_*R\Gamma(U,\mathcal G)\to j_!\mathcal G[1].
\tag{CGS8.3}
\] There are no derived Hom maps from \(i_*A\) to \(Rj_*\mathcal G\):
the exact restriction \(j^{-1}\) sends \(i_*A\) to zero, and its right
adjoint \(j_*\) sends injective modules to injectives. Thus the
corresponding Hom complex is identically zero. Likewise \(i^{-1}\) to
the full stalk ring is exact, so its right adjoint \(i_*\) preserves
injectives and is fully faithful. Applying derived Hom to CGS8.3 proves
the first isomorphism. The complementary idempotents in \(R_0\)
eliminate the \((1-\epsilon)A\) factor, proving the second.

In degree one this recovers CGS4.4 with the actual map \(r\). In higher
degrees it retains all the open cohomology and any abelian-group
extensions in the indicated derived Hom. None is replaced by its
numerical dimension.

The same gluing classification with specified actions holds in the
equivariant category. In degree one the right side is exactly the
equivariant Hom group of the map \(r\). For higher equivariant Ext, the
derived Hom on the right must be taken in the corresponding module
category with its specified action; ordinary nonequivariant Ext is not
silently substituted for it.

\subsection{CGS9. Application to the complete universal structure
sheaf}\label{cgs9.-application-to-the-complete-universal-structure-sheaf}

The actual structure sheaf has the exact datum \[
 \mathcal R
 \longleftrightarrow
 (\mathcal O_U,\mathbb Z,\mathbb Z,\mathrm{id}_{\mathbb Z}).
\tag{CGS9.1}
\] Indeed its global sections are \(R_0=\mathbb Z\oplus\mathbb Z\) and
its restriction is first projection. The open cohomology was explicitly
calculated in SS4 by the flabby resolution \[
 0\to\mathcal O_U\to\mathcal Q
 \to\bigoplus_p i_{p*}(\mathbb Q/\mathbb Z_{(p)})\to0.
\] Its global section map is onto by the retained rational
representatives for each finite residue tuple. Thus \[
 R\Gamma(U,\mathcal O_U)=\mathbb Z[0].
\tag{CGS9.2}
\] Substitution in CGS6 gives \[
 R\Gamma_{\{\mathfrak m\}}(X,\mathcal R)
 \cong\mathbb Z[0]\oplus
 [\,\mathbb Z\xrightarrow{\mathrm{id}}\mathbb Z\,],
\tag{CGS9.3}
\] where the second summand occupies degrees zero and one. Sending its
degree-one integer to the same degree-zero integer is a contracting
homotopy. Consequently its supported cohomology is exactly \(\mathbb Z\)
in degree zero and zero in every positive degree, with the explicit
restriction lift \(a\mapsto(a,0)\). This is the actual completed
application, not a theorem assuming a vanishing obstruction.

The classification also computes a different arrow that must be
retained: \[
 0\to j_!\mathcal O_U\to\mathcal R
 \to i_*(\mathbb Z\oplus\mathbb Z)\to0.
\tag{CGS9.4}
\] Its gluing class under CGS4.4 is the map \[
 (a,b)\longmapsto a,
\] or equivalently \(\mathrm{id}_{\mathbb Z}\) on the epsilon component.
It is therefore nonzero and this extension is not split as a sheaf
extension. This does not contradict CGS9.3 or SS3's split
support/restriction sequence: \[
 0\to i_*\mathbb Z\to\mathcal R
 \to j_*\mathcal O_U\to0.
\tag{CGS9.5}
\] The exact relation is CGS4.3. In CGS9.4 one tries to extend a sheaf
supported at the point into global and open sections without creating
its nonzero arithmetic restriction; naturality forbids a section
precisely because the computed map is identity. In CGS9.5 one lifts an
already present open arithmetic section by \((a,0)\); that is the proved
section. Both are actual source-induced maps and both have now been
calculated.

For complexes of coefficients the structure-sheaf identity comparison
generalizes to the explicit object CGS6.3. Once the actual Mellin or
Connes--Consani coefficient entry
\((\mathcal G^\bullet,A_+^\bullet,A_-^\bullet,r^\bullet)\) is supplied
by its proven receiving maps, that formula retains its complete defect.
This note does not invent that entry or replace it by the structure
sheaf.

\subsection{CGS10. Provenance and the exact mathematical
output}\label{cgs10.-provenance-and-the-exact-mathematical-output}

For the general gluing equivalence, the explicit human reference is
Pierre Deligne,
\href{https://www.numdam.org/item/PMIHES_1980__52__137_0/}{\emph{La
conjecture de Weil. II}}, §3.4.8, printed p.209. The local
\href{https://www.numdam.org/item/PMIHES_1980__52__137_0/}{French
record export}, lines2375--2383, was consulted for that precise passage;
it is a transcription/record export, \textbf{not original author LaTeX}.
The paper recalls the triple consisting of \(j^*\mathcal F\),
\(i^*\mathcal F\), and its specialization map
\(i^*\mathcal F\to i^*j_*j^*\mathcal F\). Its surrounding finite-field
weight argument is not thereby imported into the present receiver.

Here that specialization target has the exact, proved form \[
 i^{-1}j_*\mathcal G
 =(j_*\mathcal G)_{\mathfrak m}
 =\Gamma(X,j_*\mathcal G)
 =\Gamma(U,\mathcal G).
\tag{CGS10.1}
\] The second equality uses the sole neighbourhood \(X\) of
\(\mathfrak m\); the last is the definition of direct image along \(j\).
The full receiving closed coefficient is
\(i^{-1}\mathcal M=A_+\oplus A_-\), and its specialization is exactly
\((a_+,a_-)\mapsto r(a_+)\). Indeed it is the original restriction from
\(X\) to \(U\), and its \(A_-\) component vanishes because
\(1-\epsilon\) restricts to zero. Thus the two inverse constructions in
CGS1 prove the literal relation between Deligne's recalled triple and
the displayed quadruple, with the receiving idempotent splitting and the
entire specialization arrow retained. The distinction between full
closed stalk and selected residue remains CGS7.

The source definitions and whole-spectrum argument originate with the
user and remain in their original notation and verbatim corpus. The
universal source receiver and its sheaf are programme calculations ST
and SS. The distinction between source scalar arithmetic, power-map
pullback and the independent idèle spectral action is proved in
ST7--ST7a, with the actual original Connes--Consani source recorded
there: \href{https://arxiv.org/abs/2602.15941v1}{\emph{On the Jacobian
of \(\overline{\operatorname{Spec}\mathbb Z}\)}}, original
\href{https://arxiv.org/abs/2602.15941v1}{\nolinkurl{Jacobian.tex}}, lines 2501--2550. Their
geometric construction is not replaced by an unspecified sheaf in this
note.

The lifting mechanism under comparison is Pierre Deligne's
\href{https://www.numdam.org/item/PMIHES_1980__52__137_0/}{\emph{La
conjecture de Weil. II}}, §3.6.1--§3.6.3. The present calculation
determines which maps and coefficient objects are needed before
numerical weight separation can be invoked: the actual closed-stalk
complex, the actual open coefficient complex, and the exact restriction
map. Their complete supported receiving object is the mapping fibre
CGS6.3.

For the source structure sheaf those entries are computed, and the
supported obstruction vanishes by the explicit identity homotopy. For
every other coefficient sheaf the remaining object is no longer an
unspecified comparison: it is exactly \(\ker r\),
\(\operatorname{coker}r\) and the retained higher open cohomology, or
their full mapping-fibre complex. These statements assign no addition,
parity, numerical weight or coordinate to primitive \(Z_1/\tau\), and
make no RH or purity assertion.
